% Generated deterministically from the audited Markdown source.
% Responsible author: Carptopus.
\documentclass[11pt]{article}
\usepackage[a4paper,margin=27mm]{geometry}
\usepackage[T1]{fontenc}
\usepackage[utf8]{inputenc}
\usepackage{lmodern}
\usepackage{microtype}
\usepackage{amsmath,amssymb,amsthm,mathtools}
\usepackage{xcolor}
\usepackage[colorlinks=true,linkcolor=blue!55!black,citecolor=blue!55!black,urlcolor=blue!65!black]{hyperref}
\usepackage{xurl}
\usepackage{fancyhdr}

\newtheorem{theorem}{Theorem}
\newtheorem{lemma}[theorem]{Lemma}
\numberwithin{equation}{section}
\pagestyle{fancy}
\fancyhf{}
\fancyhead[L]{Carptopus}
\fancyhead[R]{Thue--Morse run separation}
\fancyfoot[C]{\thepage}
\setlength{\headheight}{14pt}
\setlength{\parindent}{0pt}
\setlength{\parskip}{0.42em}
\setlength{\emergencystretch}{4em}

\title{Exact first separation in the\\Thue--Morse run-length sequence}
\author{Carptopus\\\href{mailto:carptopus@163.com}{carptopus@163.com}}
\date{Manuscript, 2 September 2026}

\hypersetup{
  pdftitle={Exact first separation in the Thue-Morse run-length sequence},
  pdfauthor={Carptopus},
  pdfsubject={Exact first separation of ratio-four kernel subsequences},
  pdfkeywords={Thue-Morse sequence, run length, Jacobsthal representation, 2-kernel, automatic sequence, first separation}
}

\begin{document}
\maketitle

\begin{center}
\fcolorbox{black!35}{black!4}{\parbox{0.88\textwidth}{\centering\small
Internally verified candidate manuscript, version 0.1-beta. The mathematical claims have passed
project-internal zero-trust proof and manuscript audits; external review is pending.}}
\end{center}

\begin{abstract}


Let $d=(d(n))_{n\geq0}$ be the adjacent-run-length sequence of the Thue--Morse word, equivalently
the fixed point beginning in $1$ of $1\mapsto121$, $2\mapsto12221$. For $q\geq2$, we determine
exactly the first position at which the two kernel subsequences $d(2^{q+2}n)$ and $d(2^qn)$
differ. Writing

\[
E_q=\begin{cases}2^{q+1}-2,&q\text{ even},\\2^{q+1}+1,&q\text{ odd},\end{cases}
\]

we prove that the first position is $\lfloor2^{E_q}/9\rfloor$. The proof combines the canonical
Jacobsthal representation with a scale-stability lemma and an extremal alternating-binary-digit
argument. We extend the result to arbitrary distances in the same $4$-scale chain and derive
double-exponential lower bounds when the odd part of the multiplier is fixed and its 2-adic
valuation grows, with an exact infinite family governed by multiplicative orders.

\end{abstract}

\section{Main results}


Let

\[
N_q=\min\{n\geq1:d(2^{q+2}n)\neq d(2^qn)\}.
\]

\begin{theorem}


For every $q\geq2$,

\[
N_q=M_q:=\left\lfloor\frac{2^{E_q}}9\right\rfloor,
\qquad
E_q=\begin{cases}2^{q+1}-2,&q\text{ even},\\2^{q+1}+1,&q\text{ odd}.\end{cases}
\tag{1}
\]

\end{theorem}

The cases $q=3,4$, giving $14563$ and $119304647$, were recorded by Allouche, Allouche and
Shallit. The contribution of Theorem 1 is the exact all-$q$ formula and its minimality proof, not
the discovery of those individual values or the eventual separation of kernel subsequences.

For $r\geq1$, define in the extended positive integers, with $\inf\varnothing=+\infty$,

\[
R_r(b)=\inf\{n\geq1:d(4^{r+1}bn)\neq d(4bn)\}.
\]

\begin{theorem}


Write $b=2^ju$, where $u$ is odd, and set $q=j+2$. Then for every $r\geq1$,

\[
R_r(2^ju)\geq\left\lceil\frac{M_q}{u}\right\rceil.
\tag{2}
\]

If $u\mid M_q$, equality holds. Moreover, a fixed odd $u$ divides $M_q$ for infinitely many
even $q$ if and only if

\[
\nu_2(\mathrm{ord}_{9u}(2))\leq1.
\tag{3}
\]

\end{theorem}

\section{Two encodings of the sequence}


Let

\[
V=\{v\geq1:\nu_2(v)\equiv0\pmod2\},\qquad F(x)=|V\cap[1,x]|,
\]

and define

\[
e(x)=3F(x)-2x,\qquad H(x)=2x-F(x)=\frac{4x-e(x)}3.
\]

Then $H$ is strictly increasing and

\[
e(2x)=-e(x),\qquad e(2x+1)=1-e(x).
\tag{4}
\]

The Thue--Morse word changes between positions $n-1$ and $n$ exactly when $\nu_2(n)$ is even.
Hence, if $v_1<v_2<\cdots$ lists $V$, then $d(k)=v_{k+1}-v_k$. A direct counting argument shows

\[
d(k)=2\quad\Longleftrightarrow\quad k\in H(V).
\tag{5}
\]

If $x=\sum_j\epsilon_j2^j$, recursion (4) gives

\[
e(x)=\sum_j(-1)^j\epsilon_j.
\tag{6}
\]

This converts separation of two values of $d$ into instability of membership in $H(V)$ under
multiplication or division by four.

The equivalent Jacobsthal encoding uses $J_i=(2^i-(-1)^i)/3$. For a binary word
$w=b_m\cdots b_1$, put $[w]_J=\sum b_iJ_i$. Lemma 3.7 of
Allouche--Allouche--Shallit gives a unique representation that starts in $1$ and ends in an even
(possibly zero) number of $1$'s; we call it canonical. Their Theorem 3.8 states, in our indexing
convention, that $d(n)=1$ exactly for $n=0$ or when a Jacobsthal representation beginning in $1$
ends in a positive even number of $0$'s or a positive even number of $1$'s; otherwise $d(n)=2$.
If $B(w)=\sum b_i2^{i-1}$ and $S(w)=\sum b_i(-1)^i$, then

\[
3[w]_J=2B(w)-S(w),\qquad [w00]_J=4[w]_J+S(w).
\tag{7}
\]

\section{Explicit witnesses}


When $q$ is even, set

\[
w_q=(10)^{2^q-2}1^q,\qquad z_q=(10)^{2^q}0^{q-2}.
\]

When $q$ is odd, set

\[
w_q=(10)^{2^q}0^{q-1},\qquad z_q=(10)^{2^q-1}0100\,1^{q-1}.
\]

Each displayed word is canonical under Lemma 3.7: for even $q$, $w_q$ ends in $q$ ones and
$z_q$ ends in zero ones; for odd $q$, $w_q$ ends in zero ones and $z_q$ ends in $q-1$ ones.
Summing the geometric progressions in (7) gives

\[
[w_q]_J=2^qM_q,\qquad[z_q]_J=2^{q+2}M_q.
\]

For even $q$, the terminal constant blocks have lengths $q$ and $q-1$; for odd $q$, they again
have lengths $q$ and $q-1$, in the opposite symbols. Thus Theorem 3.8 gives

\[
d(2^qM_q)\neq d(2^{q+2}M_q).
\]

Thus $N_q\leq M_q$.

\section{Scale stability}


\begin{lemma}


Let $v\in V$.

\begin{enumerate}
\item If $2^r\mid H(v)$, $r\geq2$, and $|e(v)|<2^r$, then $4H(v)\in H(V)$.
\item If $2^r\mid H(v)$, $r\geq4$, and $|e(v)|<2^{r-2}$, then $H(v)/4\in H(V)$.

\end{enumerate}

\end{lemma}

\begin{proof}


Write $e=e(v)$. Since $3H(v)=4v-e$ and $2^r\mid H(v)$, we have $4\mid e$.
First suppose $e>0$. Write

\[
v=3\cdot2^{r-2}a+e/4.
\]

Because $0\leq e/4<2^{r-2}$, this is a carry-free high/low binary decomposition. Put

\[
T=3\cdot2^ra+e/4=4v-3e/4.
\]

Thus $T$ is obtained from $v$ by inserting two zeroes between the two blocks. Equation (6) and
the unchanged trailing-zero count give $e(T)=e(v)$ and $T\in V$, while

\[
H(T)=\frac{4T-e(T)}3=4H(v).
\]

If $e=-4b<0$, set $k=r-2$ and write

\[
v=2^k(3a-1)+(2^k-b),\qquad
T=2^{k+2}(3a-1)+(2^{k+2}-b).
\]

For $0<b<2^k$, the identity

\[
e(2^k-b)=e(2^k-1)-e(b-1)
\]

shows that the two low blocks have the same alternating digit sum because $k$ and $k+2$ have the
same parity; both also have 2-adic valuation $\nu_2(b)$. Hence again $e(T)=e(v)$, $T\in V$, and
$H(T)=4H(v)$. For $e=0$, take $T=4v$. This proves the first assertion.

For the second assertion, $|e|/4<2^{r-4}$ leaves two genuine separating zeroes that can be
deleted. Precisely, put

\[
R=(v+3e/4)/4.
\]

If $e=4b>0$, then

\[
v=3\cdot2^{r-2}a+b,\qquad R=3\cdot2^{r-4}a+b,
\]

with $0<b<2^{r-4}$. If $e=-4b<0$, then

\[
v=2^{r-2}(3a-1)+(2^{r-2}-b),\qquad
R=2^{r-4}(3a-1)+(2^{r-4}-b).
\]

The same identity, now with exponents $r-2$ and $r-4$, gives $e(R)=e(v)$ and
$\nu_2(R)\equiv\nu_2(v)\pmod2$. For $e=0$, simply take $R=v/4$. Thus $R\in V$ in every
case and $H(v)=4H(R)$, proving the second assertion.
\end{proof}


It follows that if $2^q\mid x$ and $d(4x)\neq d(x)$, then the member side $x=H(v)$ or
$4x=H(v)$ must satisfy

\[
|e(v)|\geq2^q.
\tag{8}
\]

\section{The extremal alternating-digit argument}


Let

\[
U_m=(10)^{m-1}1=\frac{4^m-1}{3}.
\]

Then $e(U_m)=m$. Among $L$ binary positions there are only $\lceil L/2\rceil$ positive-sign
positions and $\lfloor L/2\rfloor$ negative-sign positions in (6); therefore every positive
integer with $|e(x)|\geq m$ needs at least $2m-1$ binary digits. Equality with positive sign is
attained uniquely by $U_m$.

Set $m=2^q$ and

\[
v_q=2^{s_q}U_m,\qquad
s_q=\begin{cases}q-2,&q\text{ even},\\q-1,&q\text{ odd}.\end{cases}
\tag{9}
\]

The congruence $e(v)\equiv4v\pmod{2^r}$, together with the digit-length bound, proves the
following extremal statement.

\begin{lemma}


If $v\in V$, $|e(v)|\geq2^q$, and $2^q\mid H(v)$, then $v\geq v_q$. Under the
stronger divisibility $2^{q+2}\mid H(v)$, the same minimum holds for even $q$; for odd $q$,
every candidate on that member side satisfies $H(v)/4>H(v_q)$.

\end{lemma}

\begin{proof}


We use

\[
e(v)\equiv4v\pmod {2^r}.
\tag{10}
\]

First assume $2^q\mid H(v)$ and, for a contradiction, $v<v_q$. The digit-length bound reduces the
possibilities to

\[
e(v)=m+4a\quad\text{or}\quad e(v)=-m-4a,
\tag{11}
\]

where $a\geq0$. For $q\geq3$, the available number of positive or negative digit positions gives
$4a\leq(q+1)/2<2^q$, and hence $a<2^{q-2}$; when $q=2$, only $a=0$ remains. Put
$k=q-2$.

In the positive branch with $a>0$, congruence (10) gives $v=2^kA+a$, hence

\[
e(v)=(-1)^ke(A)+e(a).
\]

The assumed bound $v<v_q$ leaves at most $2m-1$ digits in $A$ when $q$ is even and at most
$2m$ when $q$ is odd. In both cases $(-1)^ke(A)\leq m$, so

\[
e(v)\leq m+e(a)\leq m+a<m+4a,
\]

a contradiction. If $a=0$, then $2^{q-2}\mid v$. Since $v\in V$, its trailing-zero count is at
least $s_q$; deleting those even-many zeroes leaves an integer with alternating digit sum at least
$m$, hence no smaller than $U_m$. Therefore $v\geq v_q$.

In the negative branch, $a=0$ is impossible: after deleting the forced trailing zeroes, the core
has at most $2m-1$ digits and starts in a positive-sign position, so its alternating digit sum
cannot be $-m$. If $a>0$, then

\[
v=2^kA+(2^k-a),\qquad
e(2^k-a)=e(2^k-1)-e(a-1).
\]

Here $e(2^k-1)$ is zero for even $q$ and one for odd $q$. Combining this with the same bound on
$A$ gives

\[
e(v)\geq-m-a+2>-m-4a,
\]

again a contradiction. This proves the first assertion.

Now assume $2^{q+2}\mid H(v)$. For even $q$, again suppose $v<v_q$, but split off the low $q$
bits in (10). In the positive branch, the low block is $2^{q-2}+a$, contributing $1+e(a)$,
while the high block has at most $2m-3$ digits and contributes at most $m-1$. Thus $a>0$
again gives $e(v)<m+4a$. When $a=0$, equality forces the high block to be $U_{m-1}$, and

\[
2^qU_{m-1}+2^{q-2}=2^{q-2}U_m=v_q.
\]

In the negative branch, the low block is $3\cdot2^{q-2}-a$. Its contribution is zero for
$a=0$ and at least $-a$ for $a>0$, while the high block contributes at least $-(m-2)$.
Neither case reaches $-m-4a$. Hence the even-$q$ minimum is still $v_q$.

Finally let $q$ be odd and suppose $v\leq4v_q$. Then $v$ has at most $2m+q$ digits, so

\[
|e(v)|\leq m+(q+1)/2.
\]

Together with (11), this implies $a<2^{q-2}$. In the positive branch, the low $q$ bits are
$2^{q-2}+a$ and contribute $-1+e(a)$, while the high block contributes at most $m$; hence

\[
e(v)\leq m-1+e(a)<m+4a.
\]

In the negative branch, $a=0$ would force $\nu_2(v)=q-2$, which is odd and contradicts
$v\in V$. For $a>0$,

\[
e(3\cdot2^{q-2}-a)=2-e(a-1)\geq3-a,
\]

so $e(v)\geq-m+3-a>-m-4a$. Thus $v>4v_q$. Since $H$ is strictly increasing and
$H(4v_q)=4H(v_q)+m$, we obtain $H(v)/4>H(v_q)$. This proves Lemma 4.
\end{proof}


We now spell out the four member-side cases needed for the lower bound. Put $x=2^qn$ and suppose
$d(4x)\neq d(x)$. By (5), exactly one side belongs to $H(V)$.

\begin{itemize}
\item If $x=H(v)$ and $q$ is odd, Lemma 4 gives
$x\geq H(v_q)=2^qM_q$.
\item If $x=H(v)$ and $q$ is even, it gives
$x\geq H(v_q)=2^{q+2}M_q>2^qM_q$.
\item If $4x=H(v)$ and $q$ is even, the strong-divisibility part gives
$x=H(v)/4\geq H(v_q)/4=2^qM_q$.
\item If $4x=H(v)$ and $q$ is odd, it gives
$x=H(v)/4>H(v_q)=2^qM_q$.

\end{itemize}

Consequently $n\geq M_q$ in every case. The witnesses in Section 3 give the reverse inequality,
proving Theorem 1. Explicitly,

\[
M_q=\begin{cases}
(4^{2^q-1}-1)/9,&q\text{ even},\\
(2\cdot4^{2^q}-5)/9,&q\text{ odd}.
\end{cases}
\]

\section{Arbitrary distances and multipliers}


The values $M_q,M_{q+2},M_{q+4},\ldots$ are strictly increasing. For $n<M_q$, no adjacent
pair in

\[
d(2^qn),d(2^{q+2}n),\ldots,d(2^{q+2r}n)
\]

has separated. At $n=M_q$, the first pair separates while all later adjacent pairs remain equal.
Therefore

\[
\min\{n\geq1:d(2^{q+2r}n)\neq d(2^qn)\}=M_q
\]

for every $r\geq1$. Substituting $b=2^ju$ proves (2), and $u\mid M_q$ gives the exact witness
$n=M_q/u$.

For even $q$, $M_q=(2^{E_q}-1)/9$. If $L=\mathrm{ord}_{9u}(2)$, then

\[
u\mid M_q\quad\Longleftrightarrow\quad L\mid2(2^q-1).
\]

This has infinitely many even solutions $q$ exactly when the 2-adic valuation of $L$ is at most
one. Indeed, write $L=2^as$ with $s$ odd. Necessity follows because
$\nu_2(2(2^q-1))=1$. Conversely, if $a\leq1$, choose $q$ divisible by
$\operatorname{lcm}(2,\operatorname{ord}_s(2))$; then $q$ is even and
$s\mid2^q-1$, so $L\mid2(2^q-1)$. When $s=1$, any even $q$ works. Taking arbitrary positive
multiples of the chosen $q$ gives infinitely many solutions and proves (3).

\section{Scope}


Theorems 1 and 2 solve an infinite ratio-four subfamily of the general first-separation question.
They do not determine the exact answer when $u\nmid M_q$, nor do they solve the problem for two
arbitrary multipliers. In the nondivisibility case, the natural finite-state descriptions retain all
residue classes modulo $u$, and no parameter-independent second-minimum theorem is presently known.

\begin{thebibliography}{9}
\footnotesize
\raggedright

\bibitem{aas2006}
G.~Allouche, J.-P.~Allouche, and J.~Shallit,
\emph{Kolam indiens, dessins sur le sable aux \^iles Vanuatu, courbe de Sierpi\'nski et morphismes de mono\"ides},
Annales de l'Institut Fourier \textbf{56} (2006), 2115--2130.
\href{https://doi.org/10.5802/aif.2235}{doi:10.5802/aif.2235}.

\bibitem{sloane}
N.~J.~A.~Sloane,
\emph{The On-Line Encyclopedia of Integer Sequences}, sequence A015565.
\href{https://oeis.org/A015565}{OEIS A015565}.

\bibitem{shallit-open}
J.~Shallit,
\emph{Open Problems in Automata Theory and Formal Languages}, Open Problem 7.
\url{https://cs.uwaterloo.ca/~shallit/Talks/open10r.pdf}.

\end{thebibliography}

\paragraph{AI-assisted research and writing disclosure.}
Carptopus is the author responsible for the mathematical claims and the public version of this
manuscript. OpenAI Codex was used as an AI-assisted tool for literature organization, proof
development and stress testing, exact verification, adversarial auditing, and manuscript
preparation.

\end{document}
